\section{An equivariant unramified eigencomplex}\label{sec:unramified}

The first specialization will start with an unramified eigencomplex on a
smooth projective curve in characteristic zero.  The geometric existence
theorem supplies such a complex, but does not by itself supply the
semilinear action needed here.  We obtain that action by using the
particular eigencomplex attached to a spectral skyscraper.  Its spectral
component is preserved by transport, and the skyscraper can be recognized
from its endomorphisms.  A second rigidity argument then shows that every
isomorphism with its transport respects the full eigenstructure.

\begin{proposition}\label{prop:equivariant-unramified}
Let $C$ be a smooth projective connected curve of genus at least two over
an algebraically closed field $K$ of characteristic zero.  Let $f$ be a
semilinear automorphism of $C$ with an ample line bundle $\mathcal L$
and an isomorphism $f^*\mathcal L\simeq\mathcal L$.  Let $G=\SL_n$
and $\Gd=\PGL_n$.  Suppose that $\tau$ is a continuous $\Gd$-local system
on $C$, defined over a finite extension of $\Q_\ell$, with Zariski-dense
image and an isomorphism $u:f^*\tau\xrightarrow{\sim}\tau$.

There is a nonzero locally bounded constructible complex $D_\tau$ on
$\Bun_G(C)$, with nilpotent singular support, a full coherent Hecke
eigenstructure of eigenvalue $\tau$, and an isomorphism
\[
                 b:f^*D_\tau\xrightarrow{\sim}D_\tau
\]
compatible with that eigenstructure and $u$.  The object $D_\tau$ is compact
in the automorphic nilpotent category.  The eigenstructure uses the
relative Satake normalization of Section~\ref{sec:conventions} and is
compatible with unit, convolution, permutations, and every fusion
diagonal.  The isomorphisms $b$ and $u$ define compatible actions of
$\Z$ on the eigencomplex and its eigenvalue.
\end{proposition}

Here $f^*$ denotes transport on the curve and on its bundle stack, with
the coefficient field $\E$ fixed.  We prove the proposition in three
steps: identify the spectral summand and its behavior under transport,
recognize the transported object without a shift, and prove uniqueness
of the coherent eigenstructure on that object.

\subsection{The spectral skyscraper and transport of the action}

Write
\[
 \mathcal S=\Loc_{\Gd}^{\mathrm{restr}}(C),\qquad
 \mathcal D_C=\Shv_{\Nilp}(\Bun_G(C)).
\]
The first stack parametrizes local systems with restricted variation;
the second is the full automorphic category with global nilpotent
singular support.  We use the following precise characteristic-zero
input from the proof of \cite[Lemma~8.1]{CT}.  The Gaitsgory--Raskin
equivalence
\begin{equation}\label{eq:unramified-equivalence}
 \mathcal D_C\simeq\IndCoh_{\Nilp}(\mathcal S)
\end{equation}
is linear for the $\QCoh(\mathcal S)$-action, and the inclusion of
$\mathcal D_C$ in the ambient automorphic category preserves compact
objects \cite[Main Theorem~1.3.9(ii), Lemma~1.3.7, and Theorem~1.1.7]{GR}.
These statements apply to geometric $\E$-adic sheaves over an arbitrary
algebraically closed characteristic-zero field
\cite[Section~2.3]{GR}.  The spectral action agrees, through universal
evaluation, with the coherent Hecke action
\cite[Main Theorem~14.3.2 and Section~14.3.3]{AGKRRV}.
As in \cite[Lemma~8.1]{CT}, we identify the external Satake labels with
our input--output convention: if a tensor relabeling $\alpha$ is required,
the external point is $\tau\circ\alpha^{-1}$.  Evaluation at the point
denoted by $\tau$ below is therefore exactly $V\mapsto V_\tau$.

Let $i_\tau:\Spec\E\to\mathcal S$ denote this point, and put
\[
       \delta_\tau=i_{\tau,*}^{\IndCoh}\E.
\]
We recall the local description used in \cite[Lemma~8.1]{CT}, since it
will also recognize the transport of this object.  The components of
$\mathcal S$ are indexed by semisimple parameters; the component of an
irreducible parameter has only one isomorphism class of geometric points
\cite[Proposition~3.7.2 and Corollary~3.7.5]{AGKRRV}.  Density makes $\tau$
irreducible and gives
\begin{equation}\label{eq:unramified-centralizer}
 \Aut(\tau)=Z_{\Gd}(\operatorname{im}\tau)=Z(\Gd)=1.
\end{equation}
The tangent complex of $\mathcal S$ at $\tau$ is
$\RGamma(C,\ad(\tau))[1]$
\cite[Section~21.2.1]{AGKRRV}.  Its degree $-1$ cohomology is zero by
\eqref{eq:unramified-centralizer}.  The invariant nondegenerate form on
$\operatorname{Lie}(\Gd)$ and Poincar\'e duality give
\[
 H^2(C,\ad(\tau))
       \simeq H^0(C,\ad(\tau))^\vee(-1)=0.
\]
Thus the open-and-closed component $\mathcal S_\tau$ containing $\tau$
has no infinitesimal automorphisms or obstructions and is a smooth
one-point formal polydisc, with tangent space $H^1(C,\ad(\tau))$.

For this formal component we use the completion description
\[
 \IndCoh(\mathcal S_\tau)
       \simeq\IndCoh(\mathbb A^d_{\E})_{\{0\}},
 \qquad d=\dim_{\E} H^1(C,\ad(\tau)),
\]
as in \cite[Sections~21.1.7--21.1.10]{AGKRRV}; the right-hand category
consists of objects set-theoretically supported at the origin.  Its
compact objects are bounded coherent complexes with this support.
Indeed, coherent complexes with this support are compact and generate the supported
category; on the smooth affine space this is also the usual generation
by the Koszul complex at the origin.  In particular their cohomology
modules have finite length.  The skyscraper $\delta_\tau$ is nonzero and
compact.  Since the formal component is smooth, it has no spectral
obstruction directions, so $\delta_\tau$ has nilpotent spectral singular
support.  This is the ind-coherent singular-support condition, rather
than the microlocal support of a constructible skyscraper.

Choose an equivalence \eqref{eq:unramified-equivalence}, and let $D_\tau$
be the inverse image of $\delta_\tau$.  It is compact in $\mathcal D_C$ and
in the ambient automorphic category.  Compact automorphic objects are
bounded constructible on every finite-type smooth chart
\cite[Appendix~F, Section~F.2.2]{AGKRRV}, so $D_\tau$ is locally bounded
constructible.  Its full eigenstructure is the one constructed in
\cite[Lemma~8.1]{CT}: the module identities
\begin{equation}\label{eq:skyscraper-evaluation}
 A\otimes\delta_\tau
       \simeq i_{\tau,*}^{\IndCoh}(i_\tau^*A),
       \qquad A\in\QCoh(\mathcal S),
\end{equation}
applied to universal evaluation over $C^I$, give
\begin{equation}\label{eq:unramified-eigenmaps}
 \Hecke_{I,(V_i)}(D_\tau)
   \simeq D_\tau\boxtimes
            \Bigl(\boxtimes_{i\in I}(V_i)_\tau\Bigr).
\end{equation}
The unit, tensor, permutation, and fusion maps come from the same
universal evaluation system.  Spectral linearity transports all of them
together.  This proves the required geometric assertions about $D_\tau$.

The remaining issue is the action of $f$.  It is enough at this stage to
preserve the summand corresponding to $\mathcal S_\tau$; no equivariance
of the chosen equivalence \eqref{eq:unramified-equivalence} is required.

\begin{lemma}\label{lem:spectral-action-transport}
Semilinear transport by $f$ identifies the automorphic spectral action
with the action for the transported curve and its restricted stack of
local systems.  Consequently, if $f^*\tau\simeq\tau$, transport preserves
the summand of $\mathcal D_C$ corresponding to $\mathcal S_\tau$.
\end{lemma}

\begin{proof}
We compare the finite-set Hecke action data that characterize the
spectral action.  These data include derived representation labels,
leg factors $\operatorname{QLisse}(C)^{\otimes I}$, and their higher
compatibilities; $\operatorname{QLisse}(C)$ is the category of
quasi-lisse sheaves of \cite[Section~8.1.1]{AGKRRV}.  Exterior product
realizes the leg factors on $C^I$.  We must transport both the Hecke
diagrams and this quasi-lisse factorization.

First consider the geometric Hecke diagrams.  Transport by $f$ is
an exact $\E$-linear equivalence preserving nilpotent singular support.
It takes every Hecke correspondence over $C^I$ to its transported
correspondence and commutes with pullback, tensor product, convolution
pushforward, and their exchange maps.  Lemma~\ref{lem:satake-transport}
compares the Satake kernels and their tensor and fusion maps.
Dualizing complexes are transported as well, so the comparison also
applies when the action is expressed with dualizing pullback conventions.
The resulting comparison is one of the actual sheaf operations and
their structural maps, including the geometric maps in the leg categories.

To pass from ordinary representations to the derived labels, use the
construction of \cite[Sections~B.3.6--B.3.8]{GKRV}, recalled in
\cite[Section~14.2.2]{AGKRRV}: bounded-derived extension, restriction
to compact objects, and then ind-extension.  On the perverse
representation objects, the preceding comparison identifies the
Satake diagrams with all their structural maps.  The universal
property of the bounded derived category extends it exactly and
preserves the coherent finite-set diagrams.  The resulting right-lax
monoidal structure maps are isomorphisms.  Restriction to compact
objects and ind-extension therefore give the comparison for the
continuous derived action.  This uses the specified extension of the
representation action; it does not identify the whole derived
spherical sheaf category with the derived representation category.

It remains to compare the quasi-lisse factorization.  Transport
identifies $\operatorname{QLisse}(C)^{\otimes I}$ with the corresponding
category for the transported curve.  The exterior-product embedding
used in \cite[Section~14.2.5]{AGKRRV} is fully faithful, so the
comparison just constructed also identifies the Hecke action with
values in these leg categories.  This is the action developed in
\cite[Sections~14.2.5--14.2.8 and Main Theorem~14.3.2]{AGKRRV}.

We have now compared the full finite-set action data.  Transport also
identifies the universal evaluation objects on the restricted stacks.
By \cite[Theorem~8.1.4]{AGKRRV}, universal evaluation gives an
equivalence between the space of $\QCoh(\mathcal S)$-actions and the
space of these data, with their higher compatibilities.  Consequently
the conjugate of the spectral action under transport is the spectral
action of the transported curve.

The component $\mathcal S_\tau$ is carried to the component of
$f^*\tau$.  When these parameters are isomorphic, its component
idempotent is preserved.  This idempotent acts as the identity on
$D_\tau$, and hence also on $f^*D_\tau$, proving the final assertion.
\end{proof}

\subsection{Recognizing the transported object}

The preceding lemma places $f^*D_\tau$ in the same spectral component as
$D_\tau$.  We now recognize its image there.  The required elementary
criterion is useful independently of the spectral setting.

\begin{lemma}\label{lem:point-object}
Let $(R,\mathfrak m)$ be a commutative local $\E$-algebra with residue
field $\E$, and let $P$ be a nonzero bounded complex of $R$-modules whose
cohomology modules have finite length.  Suppose that
\[
 \Hom_{D(R)}(P,P)=\E,\qquad
 \Hom_{D(R)}(P,P[j])=0\quad(j<0).
\]
Then $P\simeq\E[r]$ for an integer $r$.
\end{lemma}

\begin{proof}
Let $a$ and $b$ be the least and greatest degrees with nonzero
cohomology.  Any two nonzero finite-length modules over $R$ admit a
nonzero map from the first to the second: take a quotient of the first
onto $\E$ and an inclusion of $\E$ into the socle of the second.  Thus,
if $a<b$, there is a nonzero map
$h:H^b(P)\to H^a(P)$.  The truncation maps give a composite
\[
 P\longrightarrow H^b(P)[-b]
   \xrightarrow{h[-b]}H^a(P)[-b]
   \longrightarrow P[a-b].
\]
Its map on degree $b$ cohomology is $h$, so it is a nonzero negative-degree
endomorphism.  This contradicts the hypothesis.  Hence $a=b$, and $P$ is
a shift of a single finite-length module $M$.

Every element of $\mathfrak m$ acts nilpotently on $M$.  Since
$\End_R(M)=\E$, this action is a nilpotent scalar and is therefore zero.
Thus $M$ is an $\E$-vector space, and all of its $\E$-linear
endomorphisms are $R$-linear.  The equality $\End_R(M)=\E$ forces
$\dim_{\E} M=1$.
\end{proof}

For the skyscraper, and therefore for $D_\tau$, we have
\begin{equation}\label{eq:unramified-endomorphisms}
 H^0\RHom(D_\tau,D_\tau)=\E,\qquad
 H^j\RHom(D_\tau,D_\tau)=0\quad(j<0).
\end{equation}
These are the ambient automorphic Hom groups, since $\mathcal D_C$ is a
full subcategory.  Transport preserves both these groups and compactness.
Lemma~\ref{lem:spectral-action-transport} places $f^*D_\tau$ in the
$\mathcal S_\tau$ summand.  Its image under
\eqref{eq:unramified-equivalence} is therefore a bounded coherent complex
supported at the origin of the smooth formal component.  Applying
Lemma~\ref{lem:point-object} to the local ring at that origin gives
\begin{equation}\label{eq:unramified-possible-shift}
                    f^*D_\tau\simeq D_\tau[r]
\end{equation}
for some $r\in\Z$.

We eliminate the shift on a bounded open, so no global cohomological
bound on $\Bun_G(C)$ is needed.  Use the $f$-invariant ample line bundle
$\mathcal L$ in the statement.
For each sufficiently large integer $m$, let
$\mathcal U_m\subset\Bun_{\SL_n}(C)$ be the open substack on which the
associated vector bundle $\mathcal V$ satisfies
\[
 \mathcal V\otimes\mathcal L^m\text{ is globally generated},
 \qquad H^1(C,\mathcal V\otimes\mathcal L^m)=0.
\]
These conditions are open and the opens cover the bundle stack by Serre
vanishing.  Each $\mathcal U_m$ is of finite type: rank and trivialized
determinant fix the Hilbert polynomial, the displayed conditions fix
$h^0$, and a choice of basis of sections gives the usual finite-type
Quot presentation.  Adding the determinant trivialization is also a
finite-type condition.  Because $f$ preserves $\mathcal L$, every
$\mathcal U_m$ is invariant under transport.

Choose $m$ with $D_\tau|_{\mathcal U_m}\ne0$.  On a finite-type smooth
atlas of $\mathcal U_m$, local bounded constructibility gives a finite
cohomological interval; descent gives the same boundedness on
$\mathcal U_m$.  Consequently
\[
 S_m=\{j\in\Z:\mathcal H^j(D_\tau|_{\mathcal U_m})\ne0\}
\]
is finite and nonempty.  Pullback along the automorphism of
$\mathcal U_m$ is exact and conservative, so it preserves $S_m$.
Equation~\eqref{eq:unramified-possible-shift} gives $S_m=S_m-r$, which
forces $r=0$.  We have obtained an isomorphism
$b:f^*D_\tau\simeq D_\tau$.  It remains to verify its compatibility with
the Hecke eigenmaps.

\subsection{Uniqueness of the full eigenstructure}

For $I$ a finite set and $L_1,L_2$ finite-rank lisse $\E$-sheaves on
$C^I$, there is a natural external-product formula
\begin{equation}\label{eq:unramified-hom-kunneth}
 \RHom(D_\tau\boxtimes L_1,D_\tau\boxtimes L_2)
 \simeq
 \RHom(D_\tau,D_\tau)\otimes_{\E}
       \RGamma(C^I,L_1^\vee\otimes L_2).
\end{equation}
We explain its applicability to the non-quasicompact bundle stack.
On every finite-type smooth scheme chart, the restrictions of $D_\tau$
are bounded constructible, and the usual constructible Hom--K\"unneth
formula applies.  Equivalently, one can use adjunction for the projection
to the chart, proper base change for $C^I$, and the projection formula.
The second factor in \eqref{eq:unramified-hom-kunneth} is a bounded
finite-dimensional $\E$-complex, hence perfect.  Tensoring with it
commutes with the limits computing smooth descent and restriction over
finite-type opens of the bundle stack.  The chartwise formulas therefore
give the displayed global formula.  On constructible chart objects these
are also the Hom complexes in the ambient ind-sheaf category.

Both factors on the right are concentrated in nonnegative degrees, by
\eqref{eq:unramified-endomorphisms} and the fact that the $L_i$ are
ordinary lisse sheaves.  Thus
\begin{align}
 \Hom(D_\tau\boxtimes L_1,D_\tau\boxtimes L_2)
       &\simeq\Hom(L_1,L_2),\label{eq:unramified-hom-zero}\\
 H^j\RHom(D_\tau\boxtimes L_1,D_\tau\boxtimes L_2)
       &=0\qquad(j<0).\label{eq:unramified-hom-negative}
\end{align}
The first identification sends a lisse-sheaf map $a$ to
$\id_{D_\tau}\boxtimes a$.

\begin{lemma}\label{lem:unramified-eigenstructure-unique}
The complex $D_\tau$ admits a unique full coherent Hecke eigenstructure
with eigenvalue $\tau$.  Consequently every isomorphism
$f^*D_\tau\simeq D_\tau$ is compatible with the eigenstructure and the
given isomorphism $u:f^*\tau\simeq\tau$.
\end{lemma}

\begin{proof}
Compare two full eigenstructures.  For each finite set $I$ and tuple of
representations $(V_i)$, their eigenmaps differ by an automorphism of
\[
       D_\tau\boxtimes\Bigl(\boxtimes_{i\in I}(V_i)_\tau\Bigr).
\]
By \eqref{eq:unramified-hom-zero}, this automorphism is uniquely
$\id_{D_\tau}\boxtimes a_{I,(V_i)}$, where $a_{I,(V_i)}$ is an
automorphism of the lisse factor.  Naturality of the eigenstructures
makes the one-leg maps $a_V$ natural in $V$.  Compatibility with
successive Hecke operations, on the locus of distinct legs, identifies
the maps $a_{I,(V_i)}$ with exterior products of the one-leg maps.
Equality on that dense open determines maps of lisse sheaves on $C^I$.
Fusion along the diagonal then gives
\[
          a_{V\otimes W}=a_V\otimes a_W,
          \qquad a_{\E}=\id.
\]
Thus the $a_V$ form a tensor automorphism of the local-system functor
$V\mapsto V_\tau$.  Its group is the centralizer of the monodromy of
$\tau$ in $\Gd$, which is trivial by
\eqref{eq:unramified-centralizer}.  All $a_V$, and hence all
$a_{I,(V_i)}$, are identities.  The two systems of eigenmaps agree.

This also proves uniqueness at the level of coherent structures.
Equation~\eqref{eq:unramified-hom-negative} says that the mapping spaces
between the displayed targets have no positive homotopy groups.
The eigenmaps identify the successive Hecke transforms in the
compatibility diagrams with targets of the same form.  Hence there is
no additional choice of higher homotopies in the unit, convolution,
permutation, or fusion compatibilities.  The resulting natural
transformation on ordinary representation labels
has a unique exact continuous extension by the universal properties of
the stable and ind constructions defining the spectral action.

Now transport the eigenstructure of $D_\tau$ by $f$, use $u$ on its
eigenvalue, and use any isomorphism $b:f^*D_\tau\simeq D_\tau$ on the
automorphic factor.  Lemma~\ref{lem:satake-transport} makes this a full
eigenstructure on $D_\tau$ with eigenvalue $\tau$.  The uniqueness just
proved identifies it with the original one, which is the claimed
compatibility of $b$.
\end{proof}

\begin{proof}[Proof of Proposition~\ref{prop:equivariant-unramified}]
The inverse image of $\delta_\tau$ constructed above is nonzero, compact,
locally bounded constructible, and carries the coherent geometric
eigenstructure \eqref{eq:unramified-eigenmaps}.  Its membership in
$\mathcal D_C$ gives nilpotent singular support.  Spectral-action transport,
Lemma~\ref{lem:point-object}, and the invariant-open argument give an
isomorphism $b:f^*D_\tau\simeq D_\tau$.
Lemma~\ref{lem:unramified-eigenstructure-unique} proves its compatibility
with $u$ and all eigenmaps.  Iterating $b$ and its inverse defines the
action of the free cyclic group, with the corresponding iterations of
$u$ on the eigenvalue.  The eigenstructure compatibilities persist under
these iterations.  This is the asserted $\Z$-equivariant eigencomplex.
\end{proof}
